\documentclass[11pt]{article}
\usepackage[T1]{fontenc}
\usepackage{amsmath,amssymb,amsthm,microtype,booktabs}
\usepackage[margin=1in]{geometry}
\newtheorem*{theorem}{Theorem}
\newtheorem*{lemma}{Lemma}
\newtheorem*{remark}{Remark}
\newcommand{\N}{\mathbb N}
\newcommand{\vthree}{\nu_3}
\begin{document}
\section{A $3$-adic construction in dimensions four and three}
\label{sec:three-four}

We use one auxiliary walk in $\mathbb Z^3$ for both results.  The
construction is governed by the parity of the ternary digits in alternating
positions.  Write
\[
 n=\sum_{j\ge 0} d_j3^j,\qquad d_j\in\{0,1,2\},
\]
and define
\[
 a_n=(-1)^{d_1+d_3+d_5+\cdots},\qquad
 b_n=(-1)^{d_0+d_2+d_4+\cdots}.
\]
Thus $(a_n,b_n)\in\{\pm1\}^2$ is one of four states.  Put
\[
 X_n=\sum_{j<n}a_j,\qquad Y_n=\sum_{j<n}b_j,
 \qquad R_n=(X_n,Y_n,n).
\]
Appending the final ternary digit gives, for $r\in\{0,1,2\}$,
\begin{equation}\label{eq:state-recursion}
 a_{3n+r}=b_n,\qquad b_{3n+r}=(-1)^r a_n.
\end{equation}
Summing the three terms in one ternary block, and then an initial segment of
the next block, yields
\begin{equation}\label{eq:prefix-recursion}
 X_{3n+r}=3Y_n+r b_n,\qquad
 Y_{3n+r}=X_n+\varepsilon_r a_n,
 \qquad (\varepsilon_0,\varepsilon_1,\varepsilon_2)=(0,1,0).
\end{equation}

The following valuation identity is the arithmetic core of the construction.
Here $\nu_3(m)$ denotes the exponent of $3$ in a nonzero integer $m$.

\begin{lemma}[Equal-state chord identity]\label{lem:three-adic-chord}
If $m<n$ and
\[
 (a_m,b_m)=(a_n,b_n),
\]
then
\begin{equation}\label{eq:three-adic-chord}
 \nu_3\!\left((X_n-X_m)^2-3(Y_n-Y_m)^2\right)=\nu_3(n-m),
\end{equation}
and the quadratic expression on the left is nonzero.
\end{lemma}

\begin{proof}
Write $m=3u+r$ and $n=3v+s$, where $r,s\in\{0,1,2\}$.  Suppose first
that $r\ne s$.  Since the endpoint states agree, the first identity in
\eqref{eq:state-recursion} gives $b_u=b_v$.  Hence
\[
 X_n-X_m\equiv(s-r)b_u\not\equiv0\pmod3
\]
by \eqref{eq:prefix-recursion}.  Therefore
\[
 (X_n-X_m)^2-3(Y_n-Y_m)^2\not\equiv0\pmod3.
\]
At the same time $3\nmid n-m$, so both sides of
\eqref{eq:three-adic-chord} are zero.

Now suppose $r=s$.  From \eqref{eq:state-recursion}, equality of the two
endpoint states implies
\[
 (a_u,b_u)=(a_v,b_v).
\]
The correction terms in \eqref{eq:prefix-recursion} cancel, and therefore
\[
 X_n-X_m=3(Y_v-Y_u),\qquad
 Y_n-Y_m=X_v-X_u.
\]
Consequently
\[
 (X_n-X_m)^2-3(Y_n-Y_m)^2
   =-3\bigl((X_v-X_u)^2-3(Y_v-Y_u)^2\bigr),
\]
while $n-m=3(v-u)$.  Thus one division of the indices by $3$ adds one to
both valuations.  Repeating until the two indices have different residues
modulo $3$ reduces the claim to the first case.  The same argument shows
that the quadratic expression never vanishes.
\end{proof}

We shall use only the state $(1,1)$.  Let
\[
 E=\{n\ge0:(a_n,b_n)=(1,1)\}.
\]

\begin{lemma}\label{lem:no-four-selected}
No four points $R_n$ with $n\in E$ are collinear.
\end{lemma}

\begin{proof}
Suppose that four such points lie on one line.  Since their third coordinates
are distinct integers, the integer points on this line have a primitive
integer direction $D=(A,B,H)$ with $H>0$.  Hence the four points can be
written
\[
 R+q_0D,\ R+q_1D,\ R+q_2D,\ R+q_3D
\]
for distinct integers $q_0,q_1,q_2,q_3$.

For every $i<j$, Lemma~\ref{lem:three-adic-chord} applied to the
corresponding pair gives
\[
 \nu_3\!\left((q_j-q_i)^2(A^2-3B^2)\right)
   =\nu_3\!\left((q_j-q_i)H\right).
\]
In particular $A^2-3B^2\ne0$, and after cancelling one copy of
$\nu_3(q_j-q_i)$ we see that all six differences $q_j-q_i$ have the same
$3$-adic valuation, say $e$.  Dividing every $q_i-q_0$ by $3^e$ would then
produce four integers that are pairwise distinct modulo $3$, which is
impossible.
\end{proof}

It remains to understand how the indices in $E$ are spaced.  Enumerate them
as
\[
 0=N_0<N_1<N_2<\cdots.
\]
For $n=9q+3u+v$, with $u,v\in\{0,1,2\}$, the digit definition gives
\begin{equation}\label{eq:mod-nine-state}
 (a_n,b_n)=\bigl((-1)^u a_q,(-1)^v b_q\bigr).
\end{equation}
Thus, according to the state of $q$, the residues modulo $9$ that belong to
$E$ are
\[
\begin{array}{c|cccc}
(a_q,b_q)&(1,1)&(1,-1)&(-1,1)&(-1,-1)\\ \hline
n-9q&0,2,6,8&1,7&3,5&4.
\end{array}
\]
When $q$ is increased by one, a ternary carry changes the parity of exactly
one of the two alternating digit sums.  Reading consecutive entries in the
above table therefore shows that
\[
 N_{k+1}-N_k\in\{2,4,6,8\}.
\]
A direct sum of $(a_j,b_j,1)$ over each of the four possible gaps gives
\begin{equation}\label{eq:return-displacements}
\begin{array}{c|cccc}
r=(N_{k+1}-N_k)/2&1&2&3&4\\ \hline
\frac12(R_{N_{k+1}}-R_{N_k})=d_r
 &(1,0,1)&(-1,1,2)&(0,-1,3)&(-2,0,4).
\end{array}
\end{equation}
This short table will produce both walks.

\begin{theorem}[Dimension four]\label{thm:dimension-four}
There is an infinite walk in $\mathbb N^4$ using only the four standard
basis vectors as steps and having no four collinear vertices.
\end{theorem}

\begin{proof}
Let $W_0=0$.  At step $k$, if $(N_{k+1}-N_k)/2=r$, put
\[
 W_{k+1}=W_k+e_r,
\]
where $e_1,\ldots,e_4$ are the standard basis vectors of $\mathbb R^4$.
Define a linear map $T:\mathbb R^4\to\mathbb R^3$ by
\[
 T(e_r)=d_r\qquad(1\le r\le4).
\]
Equation~\eqref{eq:return-displacements} and telescoping give
\[
 T(W_k)=\frac12R_{N_k}\qquad(k\ge0).
\]

If four vertices $W_{k_0},W_{k_1},W_{k_2},W_{k_3}$ were collinear, their
images under $T$ would also be collinear.  These four images are distinct,
because their third coordinates are
\[
 N_{k_0}/2<N_{k_1}/2<N_{k_2}/2<N_{k_3}/2.
\]
This contradicts Lemma~\ref{lem:no-four-selected}.
\end{proof}

\begin{theorem}[Dimension three]\label{thm:dimension-three}
There is an infinite walk in $\mathbb N^3$ using only the three standard
basis vectors as steps such that every affine line contains at most six
vertices.
\end{theorem}

\begin{proof}
The first three vectors in \eqref{eq:return-displacements} are linearly
independent; indeed
\[
 \det(d_1,d_2,d_3)=6.
\]
They also satisfy
\[
 d_1+d_4=d_2+d_3.
\]
Let $e_0,e_1,e_2$ denote the standard basis of $\mathbb R^3$.  There is a
unique linear map $T:\mathbb R^3\to\mathbb R^3$ with
\[
 T(e_0)=d_1,\qquad
 T(e_1)=d_2,\qquad
 T(e_0+e_2)=d_3.
\]
Because both triples are linearly independent, $T$ is invertible.  The
relation above then gives
\[
 T(e_1+e_2)=d_2+d_3-d_1=d_4.
\]

Accordingly, if at step $k$ the return type is $r$, let the corresponding
macro-step be
\[
 s_1=e_0,\qquad s_2=e_1,\qquad
 s_3=e_0+e_2,\qquad s_4=e_1+e_2.
\]
Starting from $S_0=0$ and setting $S_{k+1}=S_k+s_r$, we obtain by
\eqref{eq:return-displacements}
\[
 T(S_k)=\frac12R_{N_k}\qquad(k\ge0).
\]
Since $T$ is invertible, Lemma~\ref{lem:no-four-selected} implies that no
four boundary vertices $S_k$ are collinear.

Finally subdivide the two diagonal macro-steps by
\[
 e_0+e_2=(e_2)+(e_0),\qquad
 e_1+e_2=(e_2)+(e_1).
\]
The resulting walk uses only $e_0,e_1,e_2$.  If
\(\mathcal S=\{S_k:k\ge0\}\), then every vertex of the subdivided walk is
either a macro-vertex $S_k$ or the single intermediate point $S_k+e_2$.
Hence its vertex set is contained in
\[
 \mathcal S\cup(\mathcal S+e_2).
\]
Every affine line contains at most three points of $\mathcal S$, because no
four macro-vertices are collinear.  Translation preserves collinearity, so
the same is true for $\mathcal S+e_2$.  Therefore every line contains at
most six vertices of the full unit-step walk.
\end{proof}

\begin{remark}\label{rem:d3-sharp-for-construction}
The number six is attained by this particular construction: the vertices at
times
\[
 64,70,82,88,100,106
\]
are
\[
 (38,13,13),\ (41,15,14),\ (47,19,16),\
 (50,21,17),\ (56,25,19),\ (59,27,20),
\]
and they lie on a line with primitive direction $(3,2,1)$.  Thus improving
the dimension-three bound to five points per line requires a genuinely
different construction, not merely a sharper estimate in the last paragraph.
\end{remark}

\end{document}
